%%%%%%%%%%  scan idx 610  |  folio 597  %%%%%%%%%%
\uline{plat} \(S'\to V\) \uline{tels que toute composante irréductible} \(X'_i\) \uline{de} \(X_{S'}\)\uline{, munie de la structure réduite induite, soit intègre sur} \(S'\).

En effet il existe une extension finie \(L\) du corps des fonctions \(k(S)\) telle que toute composante irréductible de \((X_L)_{\text{réd}}\) soit géométriquement intègre (d'après EGA IV 4.6.8). Soit \(S_1\) un \(S\)-schéma intègre tel que \(k(S_1)=L\) et qu'il existe un ouvert non vide \(V_1\) de \(S\) et un morphisme fini plat \(h:S_1\to V_1\). Il existe un ouvert non vide \(S_2\) de \(S_1\) tel que toute composante irréductible de \(X_{S_2}\), munie de la structure réduite induite, soit intègre sur \(S_2\) (EGA IV 9.7.7). Prenons \(V=V_1-h(S_1-S_2)\) et \(S'=h^{-1}(V)\).

Démontrons maintenant 1.2. Soient \(h:S'\to V_1\) et \(X'_i\) comme dans 1.3. Pour tout \(i\), soit \(f_i:Z_i\to X'_i\) une présentation de Chow avec \(Z_i\) intègre et \(f_i\) birationnel. Quitte à restreindre \(V\) et \(S'\), on peut supposer \(Z_i\) intègre sur \(S'\) (EGA IV 9.7.7). Soit \(f:Z=\coprod_i Z_i\longrightarrow X'=X_{S'}\) le morphisme canonique. Alors \(\underline{\Pic}_{Z/S'}=\prod_i\underline{\Pic}_{Z_i/S'}\) est représentable par un schéma en groupes somme d'ouverts quasi-projectifs de type fini.

Enfin, 1.1 dit que quitte à restreindre \(V\) et \(S'\), le morphisme \(f^*:\underline{\Pic}_{X'/S'}\longrightarrow \underline{\Pic}_{Z/S'}\) est représentable par un morphisme quasi-affine. Par suite \(\Pic_{X'/S'}\) est représentable par un schéma en groupes, somme d'ouverts quasi-projectifs de type fini~; grâce à la descente plate finie (SGA 1 VIII 7.6 et EGA II 6.6.1), \(\Pic_{X/S}|V\) l'est aussi.

\textbf{Lemme 1.4.} \uline{Soient} \(k\) \uline{un corps,} \(P\) \uline{et} \(Q\) \uline{deux} \(k\)\uline{-schémas en groupes localement de type fini et soit} \(f:P\to Q\) \uline{un morphisme quasi-affine. Alors} \(f\) \uline{est affine.}

En effet soit \(R=\Ker(f)\)~; \(R\) est un \(k\)-schéma en groupes de type fini qui est quasi-affine~; en vertu de (SGA 3 \(\mathrm{VI}_{\mathrm B}\) 11.11), \(R\) est alors affine.
